\documentclass[11pt,a4paper]{article}

\usepackage[utf8]{inputenc} % allow utf-8 input
\usepackage[T1]{fontenc}    % use 8-bit T1 fonts
\usepackage{hyperref}       % hyperlinks
\usepackage{url}            % simple URL typesetting
\usepackage{booktabs}       % professional-quality tables
\usepackage{amsfonts}       % blackboard math symbols
\usepackage{nicefrac}       % compact symbols for 1/2, etc.
\usepackage{microtype}      % microtypography
\usepackage{xcolor}         % colors
\usepackage{graphicx}       % graphics
\usepackage[normalem]{ulem} % strikethrough

\begin{document}

\section{Code-detection fixture}

\subsection{Case A - normal prose}

This is an ordinary paragraph of prose describing how the parser works. It uses the default body font and must remain a plain text item.

\subsection{Case B - prose with one inline monospaced word}

Call the printf function to print formatted output to standard out.

\subsection{Case C - paragraph in 'Source Code' style}

\begin{verbatim}
import sys
print(sys.argv)
\end{verbatim}

\subsection{Case D - fully monospaced paragraph, no code style}

\begin{verbatim}
SELECT * FROM users WHERE active = 1;
\end{verbatim}

\subsection{Case E - multi-line monospaced block}

\begin{verbatim}
def fib(n):
    a, b = 0, 1
    for _ in range(n):
        a, b = b, a + b
    return a
\end{verbatim}

\subsection{Case F - 'Source Reference' style (substring trap)}

See the original source for details.

\subsection{Case G - 'Listing Number' style, no code chars (substring trap)}

Listing 3.2

\subsection{Case H - pure-Courier prose, no code-indicative characters}

This memo is set in a typewriter face for a vintage look and feel throughout.

\end{document}